\documentclass[11pt]{amsart}

\usepackage{amsmath}
\usepackage{amssymb}
\usepackage{amsthm}
\usepackage{url}
\usepackage{hyperref}
%\usepackage{psfig}

\begin{document}
\baselineskip=16pt
\textheight=8.5in
%\parindent=0pt 
\def\sk {\hskip .5cm}
\def\skv {\vskip .05cm}
\def\cos {\mbox{cos}}
\def\sin {\mbox{sin}}
\def\tan {\mbox{tan}}
\def\intl{\int\limits}
\def\lm{\lim\limits}
\newcommand{\frc}{\displaystyle\frac}
\def\xbf{{\mathbf x}}
\def\fbf{{\mathbf f}}
\def\gbf{{\mathbf g}}

\def\dbA{{\mathbb A}}
\def\dbB{{\mathbb B}}
\def\dbC{{\mathbb C}}
\def\dbD{{\mathbb D}}
\def\dbE{{\mathbb E}}
\def\dbF{{\mathbb F}}
\def\dbG{{\mathbb G}}
\def\dbH{{\mathbb H}}
\def\dbI{{\mathbb I}}
\def\dbJ{{\mathbb J}}
\def\dbK{{\mathbb K}}
\def\dbL{{\mathbb L}}
\def\dbM{{\mathbb M}}
\def\dbN{{\mathbb N}}
\def\dbO{{\mathbb O}}
\def\dbP{{\mathbb P}}
\def\dbQ{{\mathbb Q}}
\def\dbR{{\mathbb R}}
\def\dbS{{\mathbb S}}
\def\dbT{{\mathbb T}}
\def\dbU{{\mathbb U}}
\def\dbV{{\mathbb V}}
\def\dbW{{\mathbb W}}
\def\dbX{{\mathbb X}}
\def\dbY{{\mathbb Y}}
\def\dbZ{{\mathbb Z}}

\def\la{{\langle}}
\def\ra{{\rangle}}
\def\summ{{\sum\limits}}
\def\eps{{\varepsilon}}
\def\lam{{\lambda}}
\def\tr{{\rm tr}}
\def\char{{\rm char}}
\def\Mat{{\rm Mat}}
\def\Func{{\rm Func}}
\def\Span{{\rm Span}}
\def\Ker{{\rm Ker}}
\def\Im{{\rm Im}}
\def\rk{{\rm rk}}
\def\phi{{\varphi}}
\def\row{{\rm row}}
\def\col{{\rm col}}
\def\Rowspace{{\rm Rowspace}}
\def\null{{\rm null}}


\bf\centerline{Homework \#6. Due on Thursday, Oct 15th, 11:59pm on Canvas}\rm
\vskip .1cm

\bf{Reading for this homework: }\rm Sections 3.1 and parts of 3.4.

\bf{Plan for next week (Oct 12, 14): }\rm  3.2 (the rank of a matrix), 4.2 (determinants of order $n$), possibly start 4.3 (properties
of determinants). Note: my plan is to define the determinants by the Leibniz formula
\skv
\centerline{\url{https://en.wikipedia.org/wiki/Leibniz_formula_for_determinants}}
\skv
\noindent 
instead of the cofactor expansion used in the book.

\bf\centerline{Problems: }\rm
\skv
For the first 2 problems do not assume anything about the rank except the following definitions:

\begin{itemize}
\item[(1)] Let $T:V\to W$ be a linear map where $V$ and $W$ are both finite-dimensional. The rank of $T$, denoted by $\rk(T)$, is defined
by $$\rk(T)=\dim(\Im(T))=\dim(T(V)).$$
\item[(2)] If $F$ is a field and $A\in \Mat_{m\times n}(F)$, we define its rank $\rk(A)$ by $\rk(A)=\rk(L_A)$ where
$L_A:F^n\to F^m$ is the map of left multiplication by $A$, that is, $L_A(v)=Av$ for all $v\in F^n$.
\end{itemize}

\skv
{\bf Problem 1: } Let $T:V\to W$ and $S:V\to W$ be linear maps.
\begin{itemize}
\item[(a)] Prove that $\Im(T+S)\subseteq \Im(T)+\Im(S)$ and give an example showing that the equality may not hold.
\item[(b)] Now assume that $V$ and $W$
are finite-dimensional. Prove that $$\rk(T+S)\leq \rk(T)+\rk(S).$$
\end{itemize}
\skv
{\bf Problem 2: } Let $T:V\to W$ be a linear map and $Z$ be a subspace of $W$.
\begin{itemize}
\item[(a)] Prove that $T^{-1}(Z)$ is a subspace of $V$. Here $T^{-1}(Z)$ is the preimage of $Z$,
that is, $T^{-1}(Z)=\{v\in V: T(v)\in Z\}$.
\item[(b)] Assume that $V$ and $W$ are finite-dimensional. Prove that $$\dim(T^{-1}(Z))\leq \dim (Z)+\dim(\Ker(T))$$
%and give an example where $T$ is surjective but the above inequality is strict (that is, not an equality).
{\bf Hint:} Apply the rank-nullity theorem to a suitable map obtained from $T$ by changing the domain.
\end{itemize}
\skv
{\bf Problem 3: } In parts (a) and (b) let $U,V,W$ be finite-dimensional vector spaces
and $T:U\to V$ and $S:V\to W$ linear maps. 
\begin{itemize}
\item[(a)] Prove that $$\null(ST)\leq \null(S)+\null(T)$$
where by definition $\null(R)=\dim(\Ker(R))$ for a linear map $R$. {\bf Hint:} You can express $\Ker(ST)$
in a suitable way using preimages, so that the result would follows immediately from Problem 2.
\item[(b)] Prove that $\rk(ST)\geq \rk(S)+\rk(T)-\dim(V)$.
\item[(c)] State and prove the analogue of (b) dealing with ranks of matrices.
\end{itemize}
\skv
{\bf Problem 4:} Let $F$ be a field and $m,n\in\dbN$. Recall from class that 
\begin{itemize}
\item[(i)] $\mathcal E_{i,j}(\lam)$
is the elementary row operation on the set $\Mat_{m\times n}(F)$ which
adds the $j^{\rm th}$ row multiplied by $\lam$ to the $i^{\rm th}$ row and does not change any other row
(symbolically $\row_i\mapsto \row_i+\lam\, \row_j$).
\item[(ii)] $\mathcal E'_{i,j}(\lam)$ is the elementary column operation which  adds the $i^{\rm th}$ column multiplied by 
$\lam$ to the $j^{\rm th}$ column and does not change any other column
($\col_j\mapsto \col_j+\lam\, \col_i$).
\item[(iii)] $E_{ij}(\lam)$ is the elementary (square) matrix (of suitable size) which has $1$'s on the diagonal,
$\lam$ in the $(i,j)$-position and $0$ everywhere else. Thus, $E_{ij}(\lam)=I_k+\lam e_{i,j}$ where $k$ is the size of 
the matrix and $I_k$ is the identity $k\times k$ matrix.
\end{itemize}
Prove (part of) Lemma~11.3 from class: for any $A\in \Mat_{m\times n}(F)$ we have
\begin{itemize}
\item[(a)] $(\mathcal E_{i,j}(\lam))(A)=E_{i,j}(\lam)A$ and
\item[(b)] $(\mathcal E'_{i,j}(\lam))(A)=A E_{i,j}(\lam)$,
\end{itemize}
that is, applying $E_{i,j}(\lam)$ (resp. $\mathcal E'_{i,j}(\lam)$) is the same as multiplying by
$E_{i,j}(\lam)$ on the left (resp. on the right).
\skv
{\bf Note:} The notation $E_{i,j}(\lam)$ is slightly ambiguous as it does not specify the size of the matrix. In particular,
if $m\neq n$, then $E_{i,j}(\lam)$ denotes matrices of different sizes in (a) and (b) above.
\skv
{\bf Problem 5:} Recall two other types of elementary transformations introduced in class: $\mathcal F_{i,j}$ (swappting the
$i^{\rm th}$ and $j^{\rm th}$ rows) and $\mathcal D_i(\lam)$ (multiplying the $i^{\rm th}$ row by some nonzero $\lam \in F$).
Prove that $\mathcal F_{i,j}$ can expressed as a composition of 3 operations of the form $\mathcal E_{k,l} (\lam)$
and 1 operation of the form $\mathcal D_k(\lam)$ (in a suitable order and for suitable $k,l$ and $\lam$).
\skv
{\bf Problem 6:} In Lecture~11 we explained why every matrix $A$ can be reduced to a matrix $R$ in RREF (reduced row echelon form)
using elementary row operations. The ultimate goal of this problem is to prove that such a matrix $R$ is unique. 
Given $A\in \Mat_{m\times n}(F)$, we define its rowspace $\Rowspace(A)$ to be the subspace of $F^n$ spanned by the rows of $A$.
\begin{itemize}
\item[(1)] Suppose that $A,B\in \Mat_{m\times n}(F)$ are row equivalent (that is, can be obtained from each other by a sequence
of elementary row operations). Prove that $\Rowspace(A)=\Rowspace(B)$. {\bf Note:} The converse of (1) is also true, but 
requires more work to prove.
\end{itemize}
Now let $R\in\Mat_{m\times n}(F)$ be a matrix in RREF. Our goal is to show that $R$ can be recovered from its rowspace $\Rowspace(R)$. 
Together with (1) this implies that any $A$ is row equivalent to a UNIQUE matrix in RREF.
\begin{itemize}
\item[(2)] Let $e_1,\ldots e_n$ be the standard basis of $F^n$ (which is viewed here as the space which naturally contains $\Rowspace(R)$).
For each $0\leq k\leq n$ let $V_k=\Span(e_1,\ldots, e_k)$ be the span of the first $k$ elements of the standard basis (in particular, $V_0=\{0\}$ and $V_m=F^m$). 

Now take $1\leq k\leq n$.
Prove that $V_k\neq V_{k-1}$ $\iff$ $R$ has a pivot in the $k^{\rm th}$ column. Thus, the set of columns of $R$
which contain pivots is determined by $\Rowspace(R)$. 
\item[(3)] Let $i_1<i_2<\ldots<i_p$ be the indices of columns of $R$ which contain pivots. Note that since $R$ is in RREF, this means
that the first $p$ rows of $R$ are nonzero, and the remaining rows are zero.

Let $T_0: F^n\to F^p$ be the map
which simply removes the $j^{\rm th}$ component for each $j\not\in\{i_1,\ldots, i_p\}$ (and preserves the order of the remaining components).
For example, if $n=5$, $p=2$, $i_1=2$ and $i_2=4$, we have $T_0((a_1,a_2,a_3,a_4,a_5))=(a_2,a_4)$. 

Let $T:\Rowspace(A)\to F^p$ be the restriction of $T_0$ to $\Rowspace(R)$. Prove that $T$ is an isomorphism and
for all $1\leq i\leq p$ we have $T(\row_i(R))=\eps_i$, where $\{\eps_1,\ldots, \eps_p\}$ is the standard basis of $F^p$. Thus,
$\row_i(R)$ is the unique element of $\Rowspace(R)$ which is mapped to $\eps_i$ under $T$. Since the indices $i_1,\ldots, i_p$ (and hence
also the number $p$) are determined by $\Rowspace(R)$ by (2) and all rows of $R$ below the $p^{\rm th}$ row are zero, it follows that each row of $R$ can be recovered from $\Rowspace(R)$, as desired.
\end{itemize}

\skv
{\bf Problem 7:} Let $A=\begin{pmatrix} 3 & -2 \\ 2 & -1\end{pmatrix}$. Reduce $A$ to the identity matrix only using operations of the form
$\mathcal E_{12}(\lam)$ and $\mathcal E_{21}(\lam)$ for some integer $\lam$. Then use Problem~4 to explicitly write $A$ as a product of matrices of the form $E_{12}(\lam)$ and $E_{21}(\lam)$.
\skv

{\bf Problem 8:} Let $n\in\dbN$ and $F$ a field.
\begin{itemize}
\item[(a)] Let $A\in \Mat_n(F)$ be an invertible matrix. Prove that the transposed matrix $A^t$ is also invertible and 
$(A^t)^{-1}=(A^{-1})^t$.
\item[(b)] Let $A\in \Mat_n(F)$ be such that $A^k=0$ for some $k\in\dbN$ (such matrices are called \emph{nilpotent}). Prove that
the matrix $B=I_n+A$ is invertible (as before, $I_n$ is the identity matrix) and find an explicit formula for the inverse of $A$.
\end{itemize}
Do not use any facts about inverses except the definition. In (b) you may want to start with $k=2$ (possibly followed by $k=3$) before
figuring out the general formula.
\end{document}

